% ch14-systems.tex -- File systems and persistence
\chapter{File systems and persistence}
\label{ch:systems14}
\index{inode}\index{journal}\index{fsync}\index{file system}

Think of a huge library. Without a catalogue you could only find a
book by walking every shelf. A file system is the catalogue plus the
shelving rules for a disk: it turns a flat sea of numbered blocks into
named files you can find, grow and delete. It makes three promises.
\emph{Names}: a path you can remember. \emph{Durability}: bytes that
survive a power cut. \emph{Crash behaviour}: no file left half
written. The first is easy, the other two are expensive, and this
chapter pays for them on real hardware. The opening surprise is
honest: on the author's Mac our sequential read ran at 17\,GB/s and
never touched the disk at all. It was reading the \emph{buffer
cache}, the RAM the kernel keeps copies of recent file data in, which
is the layer this whole chapter exists to explain.

\begin{wins}[What you will be able to do]
After this chapter you can explain what a path lookup does, tell a
file's name from its inode and show with a program that unlinking an
open file does not destroy it. You can compute the largest file an
indexed inode can address and the metadata a 1\,GiB file needs. You
can measure on your own machine how much of a file read comes from RAM
and how much from the device, and what \texttt{fsync} costs as a
function of write size. And you can walk a write-ahead journal through
every possible crash point and say what recovery does at each.
\end{wins}

\begin{rulebox}[Before you start]
A disk (or SSD) stores fixed-size \emph{blocks}, typically 4\,KB, and
knows nothing about files or names; the kernel builds those on top.
Recall from Book~7 ch.~\ref{ch:systems09} (the memory hierarchy) that
each level is a cache for the slower level beneath it, and from
ch.~\ref{ch:systems11} that the kernel manages memory in 4\,KB pages;
the buffer cache is that idea applied to disk blocks. Book~2 ch.~4
(sets, relations and functions) is where a mapping such as
name$\to$number was first made precise: a directory is exactly such a
function.

\emph{A short C refresher on files.} \texttt{int fd = open(path,
flags, mode)} returns a small integer, a \emph{file descriptor}, or
$-1$ on failure; \texttt{read(fd, buf, n)} and \texttt{write(fd, buf,
n)} return how many bytes they moved; \texttt{close(fd)} releases
the descriptor. \texttt{struct stat st; stat(path, \&st)} fills in a
record whose fields \texttt{st.st\_ino} (inode number),
\texttt{st.st\_nlink} (number of names) and \texttt{st.st\_size} you
will see below. \texttt{printf("\%lu", x)} prints an unsigned long.
Numbers marked \emph{Mac} are from the author's earlier laptop
(APFS); numbers marked \emph{Linux} are from a 4-vCPU container with
ext4 on a virtual disk and are kept separate.
\end{rulebox}

\begin{principle}[Think like this]
Every persistence mechanism is an answer to ``what if the power fails
\emph{right now}?'' The habit of mind is to write the steps of an
update as a numbered list and ask, for each gap between two steps,
what a person would find on the disk if the machine stopped there. We
do exactly that in the journaling section.
\end{principle}

\section{Names, inodes, and the API}
\label{sec:inode14}

Why not store a file's name and bytes together? Because then
renaming or linking a file would mean moving or copying its contents.
UNIX separates the two: the \emph{name} lives in a directory, the
\emph{contents and bookkeeping} live in an \emph{inode}.
A directory is a table of \texttt{name $\to$ inode~number}; the
inode is the file's real identity --- mode, size, and the block
numbers holding its bytes. \texttt{open()} walks the path, resolves
the inode, and returns a descriptor; \texttt{read}/\texttt{write}/
\texttt{fsync}/\texttt{close} do the rest. The real inode walker
\texttt{ch14\_inodes.py} on this book's own code tree:
(Mac) it printed lines such as
\texttt{ino=2484919 mode=0o100644 blocks=8} for \texttt{ch02\_bytes.out} ---
128 files, 758\,KB payload, one line each. On the Linux container,
over this chapter's code directory, it printed three sample lines and
a total of 216 files, 1161\,KB.

\begin{figure}[htbp]
\centering
\begin{tabular}{ll}
\toprule
inode field & role \\
\midrule
mode/uid/size & metadata, $\sim$100 bytes \\
direct[12] & first 12 block numbers \\
single indirect & $\to$ 1024 more block numbers \\
double indirect & $\to$ 1024$^2$ block numbers \\
triple indirect & $\to$ 1024$^3$ block numbers \\
\bottomrule
\end{tabular}
\caption{Classic indexed allocation: a small fixed inode plus
escalating indirect blocks. Maximum file $=$
(12+1024+1024$^2$+1024$^3$) blocks $\approx$ 4\,TB at 4\,KB blocks
--- the [hard] exercise asks for the smallest working variant.}
\label{fig:inode14}
\end{figure}


\begin{example}[Max file size, computed]
\emph{Why indirection?} An inode is small and fixed in size, but
files range from 1 byte to terabytes. The trick is to put a few block
numbers in the inode and let those point to \emph{blocks of block
numbers}. With 4\,KB blocks and 4-byte block numbers one block holds
1024 numbers. The classic inode has 12 direct pointers, one
single-indirect, one double-indirect and one triple-indirect pointer:
$12+1024+1024^2+1024^3$ blocks $\approx$ 4\,TiB from an inode of
roughly 100--128 bytes. The program \texttt{ch14\_maxfile.py} does the
arithmetic step by step:

\begin{lstlisting}[basicstyle=\ttfamily\scriptsize,caption={Real output of
\texttt{python3 ch14\_maxfile.py}.}]
pointers per indirect block : 1024
direct only                 : 48 KiB
+ single indirect           : 4.05 MiB
+ double indirect           : 4.00 GiB
+ triple indirect (max)     : 4.0039 TiB
1 GiB file = 262144 blocks of 4 KiB
  block-map metadata        : 257 blocks = 1028 KiB (+ 12 direct pointers in the inode)
  one extent (start,length) : 16 bytes
\end{lstlisting}
Check the 1\,GiB line by hand: $2^{30}/2^{12}=262\,144$ blocks; the
12 direct and 1024 single-indirect pointers cover 1036 of them, the
remaining 261\,108 need $\lceil261108/1024\rceil=255$ second-level
blocks, plus one single-indirect and one double-indirect root block:
$255+1+1=257$ blocks, i.e.\ 1028\,KiB, about 0.1\% of the file. An
\emph{extent} (start block, length) describes the same contiguous
file in 16 bytes. That is why ext4 uses extents: indexed allocation
is flexible, but for big sequential files a run-length description is
far smaller.
\end{example}

\begin{example}[Names are not files: link, symlink, unlink while open]
\emph{What a beginner expects.} Deleting a file erases it, and a
``link'' is a copy. Both are wrong, and a short program shows it.
\texttt{ch14\_links.c} creates a file, gives it a second name with
\texttt{link}, makes a symbolic link, removes both names while the
file is still open, and reads from the open descriptor.

\begin{lstlisting}[basicstyle=\ttfamily\scriptsize,caption={Real output of
\texttt{ch14\_links} (Linux container; inode numbers differ on every
run).}]
created   a         ino=1630546  nlink=1 size=12
link a->b a         ino=1630546  nlink=2 size=12
          b         ino=1630546  nlink=2 size=12
symlink   s         ino=1630562  is-symlink=1 size=1 (length of the text "a")
unlinked  a         (no such name)
read via still-open fd: 12 bytes: hello inode
fstat(fd): nlink=0 (data lives until close)
\end{lstlisting}
Read it line by line. After \texttt{link}, \texttt{a} and \texttt{b}
report the \emph{same inode number} and \texttt{nlink=2}: one file,
two names. The symlink has its \emph{own} inode whose content is just
the text ``a'' (size 1). After both names are removed, the name is
gone but the open descriptor still reads the 12 bytes, and
\texttt{nlink} has dropped to 0. The kernel frees the blocks only when
the link count is zero \emph{and} no descriptor is open, which is why
you can delete a log file a running program is writing to without
crashing it.
\end{example}

\section{The layers a read() descends}
\label{sec:layers14}

\begin{figure}[htbp]
\centering
\begin{tabular}{ll}
\toprule
layer & what it turns into what \\
\midrule
VFS & open/read/write $\to$ per-fs ops \\
file system & paths $\to$ inode $\to$ block numbers \\
buffer cache & block numbers $\to$ RAM pages \\
device driver & block numbers $\to$ NVMe commands \\
\bottomrule
\end{tabular}
\caption{The descent of a \texttt{read()}: each layer renames the
same bytes --- filename to inode to blocks to commands. Cache
lines travel up; commands travel down.}
\label{fig:layers14}
\end{figure}


\section{The buffer cache, measured}
\label{sec:cache14}

\begin{figure}[htbp]
\centering
\begin{tabular}{lrl}
\toprule
pattern & MB/s & verdict \\
\midrule
sequential read & 17\,269 & RAM-speed: page cache dominates \\
random 4KB read & 5\,950 & still cached, $3\times$ slower \\
write buffered & 6\,722 & buffer cache hides the disk \\
write + fsync/1MB & 1\,896 & real flush, $3.5\times$ slower \\
\bottomrule
\end{tabular}
\caption{Real \texttt{ch14\_fs\_io.c} numbers --- and the honest
surprise is the top row: at 17\,GB/s you are reading the
\emph{page cache}, not the disk. The buffer cache is the layer this
chapter exists to explain; fsync is the syscall that pierces it.}
\label{fig:seqrand14}
\end{figure}


\begin{example}[17\,GB/s is RAM, not disk (Mac)]
\texttt{ch14\_fs\_io.c} on the Mac: sequential 17\,269\,MB/s,
far beyond what a storage device of that era delivered, because every
read hits the page cache after the build pass. Random 4\,KB reads:
5\,950\,MB/s, still cached. The folklore ``disk is slow, RAM is
fast'' is true; the surprise is that \emph{the OS has already hidden
the disk from you}, and only \texttt{fsync} gets you past it.
\end{example}

\begin{example}[Why random 4\,KB reads are slower even from RAM (Linux)]
\emph{A tempting wrong explanation} is ``random access misses the
cache''. Nothing here misses the page cache: the whole file is in RAM.
So what else differs? The call size. Reading 4\,KB at a time takes 256
times as many system calls as reading 1\,MB at a time, and each call
has a fixed cost. We added a line that reads \emph{sequentially} in
4\,KB pieces to separate the two effects:

\begin{lstlisting}[basicstyle=\ttfamily\scriptsize,caption={Real output of
\texttt{ch14\_fs\_io} (Linux container, file in the page cache).}]
sequential read: 5269.6 MB/s
sequential 4KB read: 4002.4 MB/s
random 4KB read: 3363.0 MB/s
write (no fsync): 2630.0 MB/s
write + fsync per 1MB: 267.3 MB/s
fsync per 4KB write: 0.51 ms per fsync
\end{lstlisting}
Going from 1\,MB to 4\,KB calls costs $5270\to4002$, a factor of
1.32. Making the 4\,KB reads random costs a further $4002\to3363$,
1.19. So on this machine most of the gap is the per-call overhead,
not randomness; we did not measure why the random part costs the rest
(likely less effective prefetching and cold page-table entries).
\end{example}

\begin{figure}[htbp]
\centering
\begin{tabular}{lrl}
\toprule
pattern (Linux x86\_64 container) & MB/s & what it is \\
\midrule
sequential read, 1\,MB calls & 5\,270 & page cache \\
sequential read, 4\,KB calls & 4\,002 & page cache, 256$\times$ more calls \\
random read, 4\,KB calls & 3\,363 & page cache, scattered offsets \\
buffered write, no fsync & 2\,630 & copy into RAM \\
write + fsync per 1\,MB & 267 & device flush each time (9.8$\times$ slower) \\
write + fsync per 4\,KB & $\approx$7.7 & 0.51\,ms per fsync (343$\times$ slower) \\
\midrule
buffered read, 2nd pass & 6\,313 & page cache \\
\texttt{O\_DIRECT} read & 1\,135 & bypasses the cache: the device \\
\bottomrule
\end{tabular}
\caption{Real output of \texttt{ch14\_fs\_io} and \texttt{ch14\_direct}
on the Linux container (ext4 on a virtual disk, one run each; every
number varies by tens of percent between runs). The 4\,KB-write
figure is $4096\,\text{B}/0.51\,\text{ms}$. Reading the same file
through the cache is $6313/1135\approx5.6\times$ faster than going to
the device. Compare with the Mac numbers in
Figures~\ref{fig:seqrand14} and~\ref{fig:fsync14}.}
\label{fig:linuxio14}
\end{figure}


\begin{example}[Reaching the real device: \texttt{O\_DIRECT}]
How do you measure the disk instead of the cache? Open the file with
the \texttt{O\_DIRECT} flag, which tells the kernel to skip the page
cache (buffers and sizes must be aligned to 4096). \texttt{ch14\_direct.c}
reads the same 256\,MB file three ways:

\begin{lstlisting}[basicstyle=\ttfamily\scriptsize,caption={Real output of
\texttt{ch14\_direct} (Linux container, ext4 on a virtual disk).}]
buffered read, 1st pass :     4984 MB/s
buffered read, 2nd pass :     6313 MB/s
O_DIRECT read           :     1135 MB/s
\end{lstlisting}
Even the ``first'' buffered pass is fast because the program has just
written the file, so its pages are still cached. The honest device
rate shows only with \texttt{O\_DIRECT}: 1\,135\,MB/s, so the cache
made the same read $6313/1135\approx5.6\times$ faster. Note that this
is a virtual disk whose real hardware we cannot see; the figure is the
rate this container gets, not a property of any particular SSD.
\end{example}

\section{fsync: the syscall that pierces the cache}
\label{sec:fsync14}

\begin{figure}[htbp]
\centering
\begin{tabular}{lrr}
\toprule
fsync frequency & effective write MB/s & cost \\
\midrule
never & 6\,722 & 1$\times$ \\
per 1\,MB & 1\,896 & $3.5\times$ slower \\
per 4\,KB & $\approx$110 & $60\times$ slower (0.07\,ms each) \\
\bottomrule
\end{tabular}
\caption{fsync is cheap per call but brutal per byte: at 4\,KB
granularity, the flush dominates --- write amplification in
software form. Batching is the fix; durability-per-byte is the
tax.}
\label{fig:fsync14}
\end{figure}


A \texttt{write} call returns as soon as the bytes are copied into
the kernel's cache; the data is not yet on the device, so a power cut
loses it. \texttt{fsync(fd)} says ``do not return until this file's
data is on stable storage''. That promise is what you pay for.

(Mac) Buffered writes run at 6.7\,GB/s because the kernel only copies
into RAM. Add \texttt{fsync} per megabyte and throughput drops
$3.5\times$; per 4\,KB it drops $\sim$$115\times$ (4\,KB every 0.07\,ms is $\approx$58\,MB/s against 6\,722) --- the flush
itself becomes the workload. Databases batch commits for exactly
this reason; the cost isn't the write, it's the durability
boundary.

\begin{example}[How big must a write be before fsync stops dominating?]
\emph{The naive guess} is that fsync costs a fixed amount, so any
write size is fine. We measured instead. \texttt{ch14\_fsync\_}\allowbreak\texttt{sweep.c}
writes in sizes from 4\,KB to 1\,MB, calling \texttt{fsync} after every
write:

\begin{lstlisting}[basicstyle=\ttfamily\scriptsize,caption={Real output of
\texttt{ch14\_fsync\_sweep} (Linux container, one run; a repeat
run gave 2.9 and 10.5\,ms/fsync at 1\,MB on other occasions, so this
disk is noisy).}]
write size     fsyncs     ms/fsync       MB/s
    4096 B        512        0.679        5.8
   16384 B        512        0.586       26.7
   65536 B        128        0.568      110.0
  262144 B         32        1.029      242.9
 1048576 B          8        2.374      421.2
\end{lstlisting}
Reading the table: for 4\,KB to 64\,KB the cost per fsync is about
0.57--0.68\,ms, a \emph{fixed flush cost} that does not depend on the
size. At 1\,MiB the call takes 2.374\,ms, so the data transfer part is
about $2.374-0.57=1.8$\,ms per MiB, i.e.\ $\approx$550\,MiB/s. The
fixed cost equals the transfer time when
$0.57\,\text{ms}\times550\,\text{MiB/s}\approx0.31$\,MiB: below
roughly 300\,KiB per fsync the flush dominates, above it the data
does. Throughput climbs from 5.8 to 421\,MB/s because the same fixed
cost is spread over more bytes. This is the whole argument for
\emph{group commit}: gather many small writes and flush once.
\end{example}

\section{Journaling: never half-written}
\label{sec:journal14}

\begin{figure}[htbp]
\centering
\begin{tabular}{ll}
\toprule
step & journal state \\
\midrule
1 & write TxBegin + data + metadata records to log \\
2 & write TxEnd --- the commit point \\
3 & checkpoint: copy blocks to their real locations \\
4 & free the log region \\
\bottomrule
\end{tabular}
\caption{A journal transaction. The \emph{commit} line is the
whole trick: crash before it, the transaction is ignored; crash
after it, the replay redoes the write. The file is never
half-updated --- that is crash consistency.}
\label{fig:journal14}
\end{figure}


Before the details, recall the problem. Moving 10 from account block
X to account block Y takes two block writes, and a disk can only
promise that each \emph{single} block write is atomic. If the power
fails between them, 10 units vanish. A \emph{journal} (a write-ahead
log) fixes this by first writing down the intention, in a place the
recovery code knows to look, and only then changing the real blocks.

\begin{example}[Crash at every point, with and without a journal]
\emph{The naive attempt:} write X, then write Y. \texttt{ch14\_journal\_sim.py}
models the disk as a few blocks, stops the machine after each possible
step, runs recovery, and checks the invariant $X+Y=100$ (start
$X{=}60$, $Y{=}40$, transfer 10).

\begin{lstlisting}[basicstyle=\ttfamily\scriptsize,caption={Real output of
\texttt{python3 ch14\_journal\_sim.py}.}]
  plain    crash after step 0: X=60 Y=40 sum=100 ok
  plain    crash after step 1: X=50 Y=40 sum= 90 CORRUPT
  plain    crash after step 2: X=50 Y=50 sum=100 ok
  journal  crash after step 0: X=60 Y=40 sum=100 ok
  journal  crash after step 1: X=60 Y=40 sum=100 ok
  journal  crash after step 2: X=50 Y=50 sum=100 ok
  journal  crash after step 3: X=50 Y=50 sum=100 ok
  journal  crash after step 4: X=50 Y=50 sum=100 ok
  journal  crash after step 5: X=50 Y=50 sum=100 ok
corrupt states: plain=1, journal=0
\end{lstlisting}
\begin{figure}[htbp]
\centering
\begin{tikzpicture}[font=\footnotesize, x=1cm, y=1cm]
  \tikzset{st/.style={draw=ink, thin, minimum width=2.1cm, minimum height=0.8cm, align=center, fill=accentsoft}}
  \node[st] (s1) at (0.0,1.6) {1. write\\log};
  \node[st] (s2) at (2.5,1.6) {2. write\\commit};
  \node[st,fill=goodbg] (s3) at (5.0,1.6) {3. write X\\in place};
  \node[st,fill=goodbg] (s4) at (7.5,1.6) {4. write Y\\in place};
  \node[st,fill=graybg] (s5) at (10.0,1.6) {5. clear\\log};
  \foreach \a/\b in {s1/s2,s2/s3,s3/s4,s4/s5}{\draw[->,thick] (\a) -- (\b);}
  \draw[amber, thick, dashed] (1.25,0.9) -- (1.25,2.7);
  \node[amber, anchor=north, align=center] at (1.25,0.85) {crash here:\\log has no commit\\$\to$ discard};
  \draw[good, thick, dashed] (3.75,0.9) -- (3.75,2.7);
  \node[good, anchor=north, align=center] at (6.25,0.85) {crash anywhere after the commit line:\\replay the log\\$\to$ X and Y both new};
  \node[good, anchor=south] at (3.75,2.7) {commit point};
\end{tikzpicture}
\caption{The five steps of the journaled transfer in
\texttt{ch14\_journal\_sim.py}. Before the dashed commit line a crash
means ``nothing happened''; after it a crash means ``replay it''.
There is no step at which recovery leaves $X+Y\ne100$. Without a
journal, steps 3 and 4 are the only two writes, and a crash between
them is exactly the corrupt state.}
\label{fig:crashpoints14}
\end{figure}

The plain version has exactly one bad moment, between its two writes,
and 90 is not an amount anyone ever had. The journal version writes
the new values to a log (step 1), then a \emph{commit record} (step 2),
then updates X and Y in place (3, 4), then clears the log (5). Crash
after step 1: the log has no commit, recovery discards it, nothing
happened. Crash after step 2, 3 or 4: the commit exists, recovery
\emph{replays} the logged values into X and Y; replaying is harmless
if some of it was already done, a property called idempotence. The
model assumes a single block write is atomic and does not model
torn writes within a block, which is the guarantee real disks give
and journals rely on.
\end{example}

The commit record is the whole trick: a crash before it discards
the transaction; a crash after it replays the write. Ordering the
four steps costs one extra disk write --- the journal --- and buys
atomicity across arbitrarily many blocks. xv6's \texttt{log.c}
implements exactly this; the guided-read panel shows it distilled.

\begin{figure}[htbp]
\centering\small
\fbox{\begin{minipage}{0.86\linewidth}
\textbf{Guided read --- xv6 \texttt{file.c} \& \texttt{log.c}}
(20--30 lines):\\[2pt]
\texttt{filewrite}: loop \texttt{begin\_op(); ilock(); writei();
iunlock(); end\_op();} --- a transaction per write batch.\\
\texttt{log.c}: \texttt{log\_write} appends; \texttt{commit} does
\texttt{write\_log} + \texttt{install\_trans} --- exactly the
4-step transaction of Figure~\ref{fig:journal14}, in a real kernel.\\
repo: \texttt{github.com/mit-pdos/xv6-riscv}
\end{minipage}}
\caption{The same transaction in a real kernel, distilled: xv6's
\texttt{log.c} is the journal figure with lock calls added.}
\label{fig:xv6fs14}
\end{figure}


\section{SSDs: why random writes hurt}
\label{sec:ssd14}

Flash erases in large blocks but writes in small pages --- a
4\,KB random write can trigger a read-modify-write of a whole
erase block. Sequential patterns let the FTL coalesce; random ones
force garbage collection. The shape of the penalty is the same as
ch10's line-fill argument at a different layer: granularity costs
when the device's minimum operation exceeds yours.

\section{Common mistakes}
\label{sec:mistakes14}
\begin{itemize}
\item \textbf{``The write is durable'' without fsync.} A
      \texttt{write()} returns after a RAM copy --- power loss
      loses it silently. fsync is the durability boundary.
\item \textbf{Measuring ``disk speed'' through the page cache.}
      Our own first run read 17\,GB/s (Mac) --- of RAM. A valid storage
      benchmark must outsize RAM or bypass the cache.
\item \textbf{fsync-per-write by default.} 0.07\,ms each is fine
      for 100 writes; at 4\,KB writes it's a $\sim$$115\times$ cut in the
      throughput --- batch commits.
\item \textbf{Treating directories as files' contents.} The
      directory holds \emph{names}; the inode holds the file ---
      hard links are just two names for one inode.
\item \textbf{``SSD has no seek time so order doesn't matter.''}
      Random \emph{writes} still pay erase-block amplification;
      only reads are truly seek-free.
\end{itemize}

\section{Summary}
\label{sec:summary14}
\begin{tabularx}{\linewidth}{@{}l X@{}}
\toprule
\textbf{Piece} & \textbf{Verified version} \\
\midrule
inode    & name$\to$inode$\to$blocks; 4\,TB from 100\,B \\
buffer cache & sequential reads at 17\,GB/s = RAM, measured \\
fsync    & the durability syscall; $\sim$$115\times$ tax at 4\,KB \\
journal  & TxBegin..TxEnd..checkpoint = atomic crash safety \\
SSD      & erase-block granularity punishes random writes \\
\bottomrule
\end{tabularx}

\section*{Further reading}
\label{sec:further14}
\begin{itemize}
\item OSTEP ch36--46 (the FS book-within-the-book):
      \url{https://pages.cs.wisc.edu/~remzi/OSTEP/}
\item xv6 ch8 + \texttt{file.c}/\texttt{log.c} source:
      \url{https://github.com/mit-pdos/xv6-riscv}
\item CS:APP ch10 (system-level I/O):
      \url{https://csapp.cs.cmu.edu/}
\end{itemize}

\begin{rulebox}[Watch instead]
\begin{itemize}
\item \textbf{File Systems --- MIT 6.1810 lecture} --- the xv6 FS
      walk-through that pairs with this chapter's guided read
      (verified 2026-10-11).
      \url{https://youtu.be/3b4vEyb1I4o}
\end{itemize}
\end{rulebox}

\section*{Notes}
\label{sec:notes14}
\begin{itemize}
\item The first set of throughput numbers is verbatim
      \texttt{ch14\_fs\_io.c} output on the author's Mac (APFS volume
      plus page cache); the Linux listings are from an ext4
      container on a virtual disk, labelled as such. Your numbers
      will differ --- the exercise is to compare shapes, not digits.
\item The inode walker output is real \texttt{lstat} data on this
      repo's own code directory. The journal model checks crash
      points between whole block writes only.
\end{itemize}

\section*{Exercises}
\addcontentsline{toc}{section}{Exercises}
\stepcounter{section}

\begin{exercise}[easy]
Walk a \texttt{read("/a/b/f")} down the layers: which two lookups
turn the path into an inode, and the inode into bytes?
\end{exercise}
\begin{solution}[hint]
Directory lookup per component (name$\to$ino); inode's block
map (file offset$\to$disk block).
\end{solution}

\begin{exercise}[easy]
Why does \texttt{unlink} of an open file keep working? Name the
refcounting idea.
\end{exercise}
\begin{solution}[hint]
Open descriptors pin the inode; storage frees at link \emph{and}
descriptor count zero.
\end{solution}

\begin{exercise}[easy]
Sequential write is buffered at 6.7\,GB/s on the Mac (2.6\,GB/s on the Linux container) --- where do the
bytes physically go between \texttt{write} and \texttt{fsync}?
\end{exercise}
\begin{solution}[hint]
Kernel page cache, marked dirty; the device sees them only at
flush or write-back.
\end{solution}

\begin{exercise}[medium]
Run the sequential/random read test on your own machine. Report
both numbers and name the layer you measured; how would you
\emph{force} a real disk read?
\end{exercise}
\begin{solution}[hint]
Files bigger than RAM, or \texttt{O\_DIRECT}/madvise to bypass
the cache.
\end{solution}

\begin{exercise}[medium]
A crash lands after TxEnd but before checkpoint. Reconstruct the
journal replay --- which writes happen, and why is the state
correct?
\end{exercise}
\begin{solution}[hint]
Committed transactions replay idempotently; uncommitted ones
never ran.
\end{solution}

\begin{exercise}[medium]
Design the journal for a key-value store: what goes in each
record, and what makes the commit record atomic?
\end{exercise}
\begin{solution}[hint]
Op+key+value+txid per record; the commit is a single sector
write --- atomic by hardware.
\end{solution}

\begin{exercise}[medium]
Using the Linux sweep table (flush $\approx0.57$\,ms, and
2.374\,ms for a 1\,MiB write), estimate the write size at which the
data transfer time equals the fixed flush cost, and say what the
Mac's 0.07\,ms flush implies about its crossover.
\end{exercise}
\begin{solution}[hint]
Transfer rate $\approx1\,\text{MiB}/(2.374-0.57)\,\text{ms}\approx554$
MiB/s, so crossover $\approx0.57\,\text{ms}\times554\approx0.31$\,MiB.
With the same device rate, a flush of 0.07\,ms would cross over at
about $0.07\times554\approx39$\,KiB; the device rate on the Mac was
not measured, so this is a scaling statement, not a Mac prediction.
\end{solution}

\begin{exercise}[medium]
Ext4's extents replace the block map with runs. For a 1\,GB
sequential file, compare extent vs inode+indirect metadata sizes.
\end{exercise}
\begin{solution}[hint]
One extent record (start,len) vs 256K block numbers --- orders
of magnitude.
\end{solution}

\begin{exercise}[hard]
Design the smallest inode that addresses a 1\,GB file with 4\,KB
blocks: pick your block-pointer scheme, count the metadata bytes,
and justify the minimum.
\end{exercise}
\begin{solution}[hint]
1\,GB/4\,KB = 256K blocks; one extent or a 2-level indirect tree
--- argue either, price both.
\end{solution}

\begin{exercise}[hard]
SSD FTLs do copy-on-write internally. Argue why a journaling FS
\emph{plus} a COW device can double-write the same data, and
name the classic remedy.
\end{exercise}
\begin{solution}[hint]
Journal + FTL both relocate; log-structured FS or letting the
FTL handle it (ext4 without full data journaling).
\end{solution}

\begin{exercise}[easy]
What two numbers in \texttt{lstat} output tell a hard link from
a copy?
\end{exercise}
\begin{solution}[hint]
Same inode number + link count $>$1 = one file, two names.
\end{solution}

\begin{exercise}[medium]
Compare the journal's ``write everything twice'' cost against the
alternative of checksum-only metadata ordering. When is the second
option safe?
\end{exercise}
\begin{solution}[hint]
When data ordering is already guaranteed (or data loss tolerable)
--- metadata journaling exists for this.
\end{solution}

\begin{exercise}[medium]
The walker shows \texttt{blocks=8} for a 1-byte file: 4\,KB used.
What is this called, and what does it cost a million-small-files
workload?
\end{exercise}
\begin{solution}[hint]
Internal fragmentation --- block granularity tax, as severe as
the SSD's erase-block version.
\end{solution}

\begin{exercise}[hard]
Argue whether \texttt{fsync} should return before or after the
device's internal flush, and what can still go wrong if the drive
lies about having flushed (hint: power-fail protected caches).
\end{exercise}
\begin{solution}[hint]
After flush to non-volatile media; a lying drive trades
durability for benchmark speed --- enterprise drives add
capacitors to keep the promise honest.
\end{solution}

\begin{exercise}[easy]
In the \texttt{ch14\_links} output, which two fields prove that
\texttt{a} and \texttt{b} are the same file, and which field shows
the symlink is \emph{not}?
\end{exercise}
\begin{solution}[hint]
Same \texttt{ino} and \texttt{nlink=2} for \texttt{a} and \texttt{b};
the symlink \texttt{s} has a different inode number and its size (1) is
the length of the stored path text.
\end{solution}
